\begin{lstlisting}[caption={Ch\_genererCheminOptimal}]
Ch_Chemin* Ch_genererCheminOptimal(V_Ville *v) {
    if (!v) return Ch_creerChemin();

    unsigned int nb_cases = V_obtenirNbCases(v);
    unsigned int depart = V_obtenirNumCaseDepart(v);
    
   
   // Ici on cherche tous les objectifs obligatoires  
    unsigned int *objectifs = malloc((nb_cases + 1) * sizeof(unsigned int));
    objectifs[0] = depart;  // Le départ est toujours le premier
    
    int nb_objectifs = 1;
    for (unsigned int i = 1; i <= nb_cases; i++) {
        C_Case *c = V_obtenirCase(v, i);
        if (C_estObligatoire(*c) && i != depart) {
            objectifs[nb_objectifs++] = i;
        }
    }
    
    //  Le cas où aucun objectif
    if (nb_objectifs <= 1) {
        Ch_Chemin *res = construire_chemin_final(v, (int*)&depart, 0, depart, 
                                                 V_obtenirDirDepart(v));
        free(objectifs);
        return res;
    }
    
    
    // On calcule les distances entre objectifs avec BFS
    
    int **distances = malloc(nb_objectifs * sizeof(int*));
    unsigned int ***predecesseurs = malloc(nb_objectifs * sizeof(unsigned int**));
    
    for (int i = 0; i < nb_objectifs; i++) {
        distances[i] = malloc(nb_objectifs * sizeof(int));
        predecesseurs[i] = malloc(nb_objectifs * sizeof(unsigned int*));
        
        for (int j = 0; j < nb_objectifs; j++) {
            if (i == j) {
                distances[i][j] = 0;
                predecesseurs[i][j] = NULL;
            } else {
                ResultatBFS res = bfs_chemin_court(v, objectifs[i], objectifs[j]);
                distances[i][j] = res.distance;
                predecesseurs[i][j] = res.predecesseurs;
                
                if (res.distance == INT_MAX) {
                    // libération de la mémoire
                    free(objectifs);
                    for (int k = 0; k <= i; k++) {
                        for (int l = 0; l < nb_objectifs; l++) {
                            if (predecesseurs[k][l]) free(predecesseurs[k][l]);
                        }
                        free(predecesseurs[k]);
                        free(distances[k]);
                    }
                    free(predecesseurs);
                    free(distances);
                    return Ch_creerChemin();
                }
            }
        }
    }
    
    // Ici on cherche le meilleur ordre des visites des points obligatoires
    int *sequence_temp = malloc(nb_objectifs * sizeof(int));
    int *meilleure_sequence = malloc(nb_objectifs * sizeof(int));
    bool *visites = calloc(nb_objectifs, sizeof(bool));
    int min_cout = INT_MAX;
    
    sequence_temp[0] = 0; 
    visites[0] = true;
    
    explorer_ordre_objectifs(1, nb_objectifs, sequence_temp, visites, 0,
                            distances, meilleure_sequence, &min_cout);
    
    if (min_cout == INT_MAX) {
        fprintf(stderr, "Aucune solution trouvée\n");
        free(sequence_temp);
        free(meilleure_sequence);
        free(visites);
        free(objectifs);
        for (int i = 0; i < nb_objectifs; i++) {
            for (int j = 0; j < nb_objectifs; j++) {
                if (predecesseurs[i][j]) free(predecesseurs[i][j]);
            }
            free(predecesseurs[i]);
            free(distances[i]);
        }
        free(predecesseurs);
        free(distances);
        return Ch_creerChemin();
    }
    
    // On assemble le chemin complét
    
    int taille_max = (int)(nb_cases * 2) + 2;
    int *chemin_final = malloc(taille_max * sizeof(int));
    int longueur_totale = 0;
    
    // Pour chaque segment entre objectifs
    for (int i = 0; i < nb_objectifs; i++) {
        int idx_depart = meilleure_sequence[i];
        int idx_arrivee;

        if (i == nb_objectifs - 1) {
            idx_arrivee = 0;
        } else {
        idx_arrivee = meilleure_sequence[i + 1];
        } 
        
        int *segment = malloc((nb_cases + 1) * sizeof(int));
        int longueur_segment = reconstruire_chemin(objectifs[idx_depart],
                                                   objectifs[idx_arrivee],
                                                   predecesseurs[idx_depart][idx_arrivee],
                                                   segment);
        
        if (longueur_segment < 0) {
            free(segment);
            continue;
        }
        
        // Copier le segment et permet d'éviter les doublons aux jonctions
        int debut;
        if (i == 0) {
        debut = 0;
        } else {
        debut = 1;
        }
        for (int j = debut; j <= longueur_segment; j++) {
            chemin_final[longueur_totale++] = segment[j];
        }
        
        free(segment);
    }
    
    // Construction du chemin final
    Ch_Chemin *resultat = construire_chemin_final(v, chemin_final, 
                                                  longueur_totale - 1,
                                                  depart, V_obtenirDirDepart(v));
    
    // libération de la mémoire
    free(objectifs);
    free(chemin_final);
    free(sequence_temp);
    free(meilleure_sequence);
    free(visites);
    
    for (int i = 0; i < nb_objectifs; i++) {
        for (int j = 0; j < nb_objectifs; j++) {
            if (i != j && predecesseurs[i][j] != NULL) {
                free(predecesseurs[i][j]);
            }
        }
        free(predecesseurs[i]);
        free(distances[i]);
    }
    free(predecesseurs);
    free(distances);

    return resultat;
}
\end{lstlisting}